\documentclass[11pt]{article}

\usepackage[margin=1in]{geometry}
\usepackage[T1]{fontenc}
\usepackage[utf8]{inputenc}
\usepackage{lmodern}
\usepackage{amsmath,amssymb,mathtools}
\usepackage{array}
\usepackage{booktabs}
\usepackage{longtable}
\usepackage{enumitem}
\usepackage{xcolor}
\usepackage{hyperref}
\usepackage[normalem]{ulem}
\usepackage{listings}
\usepackage{fancyhdr}
\usepackage{setspace}

\hypersetup{
  colorlinks=true,
  linkcolor=blue!50!black,
  urlcolor=blue!50!black
}

\definecolor{noteink}{RGB}{25,25,45}
\definecolor{notered}{RGB}{130,30,30}
\definecolor{codebg}{RGB}{245,242,236}
\definecolor{rulegray}{RGB}{100,100,100}

\setlength{\parskip}{0.55em}
\setlength{\parindent}{0pt}
\setlength{\marginparwidth}{1.15in}
\onehalfspacing

\pagestyle{fancy}
\fancyhf{}
\fancyhead[L]{QAOA Portfolio Optimisation Research Journal}
\fancyhead[R]{version 0.4.0 repo notebook}
\fancyfoot[C]{\thepage}

\lstset{
  basicstyle=\ttfamily\small,
  backgroundcolor=\color{codebg},
  frame=single,
  framerule=0.2pt,
  rulecolor=\color{rulegray},
  columns=fullflexible,
  keepspaces=true,
  showstringspaces=false,
  breaklines=true
}

\newcommand{\entry}[2]{\subsection*{#1 --- #2}}
\newcommand{\todo}[1]{\textcolor{notered}{\textbf{TODO:} #1}}
\newcommand{\mnote}[1]{\marginpar{\raggedright\footnotesize\color{notered} #1}}
\newcommand{\repo}[1]{\texttt{#1}}

\begin{document}

\begin{center}
{\LARGE Research Journal / Working Notebook}\\[0.4em]
{\large QAOA-based portfolio optimisation}\\[0.6em]
{\normalsize Compiled from the actual implementation path of the repository}\\
{\normalsize \repo{src/portfolio\_qaoa\_bench/} through the final repo state}
\end{center}

\hrule
\vspace{0.8em}

This is intentionally not a polished thesis chapter. It is the notebook I wish I had started with:
day-by-day enough to see how the repository actually grew, mathematical enough to be useful,
and honest enough to leave in the dead ends. I am not trying to make the project look linear.
It was not linear. It was a chain of ``I think this is right'', then a test failed, then I had
to work backwards and admit what I had not understood yet.

I have also tried to tie every major notebook claim to something that ended up in the code:
\repo{qubo.py}, \repo{simulator.py}, \repo{objective.py}, \repo{optimizers.py},
\repo{landscape.py}, \repo{pipeline.py}, \repo{reporting.py}, and later the repo-hygiene fixes
such as naming cleanup, removal of fake security modules, CI strengthening, and the generated
\repo{monolith-full/} mirror.

\section*{Week 0 --- reading, vague panic, and trying to say ``ansatz'' without bluffing}

\entry{Day 0, late evening}{What is QAOA actually doing?}

I started with Farhi et al.\ 2014 and immediately realised I had been using the word
``ansatz'' like a decoration rather than a precise thing. For this project I think the clean
way to say it is:

\[
U(\boldsymbol{\gamma},\boldsymbol{\beta})
= \prod_{\ell=1}^{p} e^{-i\beta_\ell B} e^{-i\gamma_\ell C}
\]

where \(C\) is the cost Hamiltonian and \(B\) is the mixer Hamiltonian.

My first wrong note was: ``\(U = \prod_{\ell=1}^{p} e^{-i\gamma_\ell B} e^{-i\beta_\ell C}\)''.
That order is backwards. I only caught it because I kept checking the paper and the order in the
repository sketch I wanted to build.

\mnote{ask supervisor the difference between ``variational circuit'' and ``ansatz''}

\entry{Day 1, still reading}{VQE vs QAOA, which I kept merging in my head}

My short version after too much reading:
\begin{itemize}[leftmargin=1.4em]
\item VQE is the wider template: choose a parametrised state, evaluate expectation values, optimise.
\item QAOA is the more structured special case: cost layer, mixer layer, repeated \(p\) times.
\item In this repo the structure matters because the constraint handling is not an afterthought;
the mixer choice changes whether I stay in the feasible Hamming-weight sector.
\end{itemize}

First circuit doodle in my notebook was embarrassingly generic:

\begin{lstlisting}
|0> --H----Rz(gamma)----Rx(beta)---
|0> --H----Rz(gamma)----Rx(beta)---
|0> --H----Rz(gamma)----Rx(beta)---
\end{lstlisting}

That is not the portfolio circuit. It ignores the quadratic couplings, ignores the constraint,
and ignores the fact that \(\beta\) and \(\gamma\) are not just arbitrary decorations. The more
accurate mental picture for the repo eventually became:

\begin{lstlisting}
prepare initial state (uniform OR Dicke OR feasible basis)
for each layer:
    cost:  Rz terms + Rzz terms from QUBO->Ising
    mixer: Rx on each qubit  OR  Rxx/Ryy on XY edges
measure
\end{lstlisting}

\section*{Week 1 --- QUBO construction, penalty math, and the first real feeling that the code might work}

\entry{Mon 6 Oct, 9:10am}{Penalty expansion by hand}

The budget constraint is \(\sum_i x_i = B\) for binary \(x_i \in \{0,1\}\).
The penalty term is
\[
P\left(\sum_i x_i - B\right)^2.
\]

I first wrote
\[
\left(\sum_i x_i - B\right)^2
\;=\;
\sum_i x_i^2 - 2B \sum_i x_i + B^2
\]
and then stopped because something felt too small. I had forgotten the cross terms.

The actual expansion is
\begin{align*}
\left(\sum_i x_i - B\right)^2
&=
\left(\sum_i x_i\right)^2 - 2B \sum_i x_i + B^2 \\
&=
\sum_i x_i^2 + 2\sum_{i<j} x_i x_j - 2B\sum_i x_i + B^2.
\end{align*}

Since \(x_i^2 = x_i\) for binary variables,
\begin{align*}
\left(\sum_i x_i - B\right)^2
&=
\sum_i x_i (1 - 2B) + 2\sum_{i<j} x_i x_j + B^2.
\end{align*}

So in the symmetric QUBO matrix:
\begin{itemize}[leftmargin=1.4em]
\item each diagonal entry gets \(+P(1-2B)\),
\item each off-diagonal pair contributes \(+P\) in each symmetric location,
\item the constant \(PB^2\) can be ignored for optimisation.
\end{itemize}

This is exactly what ended up in \repo{src/portfolio\_qaoa\_bench/qubo.py}:
diagonal shift \(P(1-2B)\), off-diagonal shift \(+P\), and then re-symmetrisation.

\entry{Tue 7 Oct, after lunch}{Checking one tiny \(3\times 3\) example}

I wanted one hand-computed case so I would stop treating the matrix as magic.
Suppose the base logical objective matrix is
\[
Q_{\text{base}} =
\begin{bmatrix}
-0.0940 & 0.0012 & 0.0006 \\
0.0012 & -0.1128 & 0.0009 \\
0.0006 & 0.0009 & -0.0746
\end{bmatrix}
\]
with \(B=1\) and \(P=0.2\).

Then
\[
P(1-2B) = 0.2(1-2) = -0.2
\]
and off-diagonal shift is \(+0.2\).

So the penalised QUBO becomes
\[
Q =
\begin{bmatrix}
-0.2940 & 0.2012 & 0.2006 \\
0.2012 & -0.3128 & 0.2009 \\
0.2006 & 0.2009 & -0.2746
\end{bmatrix}.
\]

That was the first time the scale clicked. The penalty dominates the couplings on purpose.
If the budget is violated, I want that violation to be expensive enough that the optimiser does
not get to pretend an infeasible portfolio is a valid low-energy answer.

\entry{Thu 9 Oct, 8:45pm}{How expensive is the exact reference really?}

The exact feasible optimum in the code does not enumerate all \(2^n\) states.
It only enumerates the feasible supports:
\[
\min_{x \in \{0,1\}^n,\; \sum_i x_i = B} x^\top Q_{\text{base}} x.
\]

That means the work is \(\binom{n}{B}\), not \(2^n\). This matters.

For example:
\[
\binom{12}{4} = \frac{12 \cdot 11 \cdot 10 \cdot 9}{4 \cdot 3 \cdot 2 \cdot 1} = 495
\]
which is completely reasonable.

But
\[
\binom{15}{7} = 6435
\]
and if I keep pretending the exact reference is always available I will eventually lie to myself
about approximation quality.

This eventually became the explicit config field
\repo{exact\_reference\_max\_assets} and the repository now uses
\repo{exact\_feasible\_energy = None} above that cutoff instead of silently faking certainty.

\section*{Week 2 --- statevector execution, Dicke states, and the XY mixer factor-of-two argument}

\entry{Mon 13 Oct, 10:20am}{Why the Dicke initial state is the right thing for the XY path}

If the XY mixer is supposed to preserve Hamming weight, then the cleanest initial state is a
uniform superposition over all budget-feasible basis states:
\[
\lvert D^n_B \rangle
=
\frac{1}{\sqrt{\binom{n}{B}}}
\sum_{\substack{x \in \{0,1\}^n \\ |x| = B}}
\lvert x \rangle.
\]

For \(n=4\), \(B=2\),
\[
\binom{4}{2} = 6
\qquad \Rightarrow \qquad
\text{amplitude} = \frac{1}{\sqrt{6}} \approx 0.408248.
\]

For \(n=6\), \(B=2\),
\[
\binom{6}{2}=15,
\qquad
\text{amplitude} = \frac{1}{\sqrt{15}} \approx 0.258199.
\]

That is exactly what the code does for small enough systems through
\repo{\_dicke\_statevector(...)} and \repo{\_prepare\_xy\_initial\_state(...)}.
If \(n\) gets larger than \repo{dicke\_init\_max\_assets}, the code falls back to a single feasible
basis state. Not as elegant, but explicit.

\entry{Tue 14 Oct, too much coffee}{Deriving the XY mixer in the \(\{|01\rangle,|10\rangle\}\) subspace}

This took me longer than it should have because I kept losing a factor of \(2\).

I wanted to understand why the statevector path uses
\[
\cos(2\beta), \qquad \sin(2\beta)
\]
inside \repo{\_apply\_xy\_layer(...)} and why the circuit path uses
\repo{rxx(2*beta)} and \repo{ryy(2*beta)}.

Start with
\[
X \otimes X \lvert 01 \rangle = \lvert 10 \rangle,
\qquad
Y \otimes Y \lvert 01 \rangle = \lvert 10 \rangle.
\]

So
\[
(X\otimes X + Y\otimes Y)\lvert 01\rangle = 2 \lvert 10 \rangle
\]
and similarly
\[
(X\otimes X + Y\otimes Y)\lvert 10\rangle = 2 \lvert 01 \rangle.
\]

Therefore in the two-dimensional subspace \(\{|01\rangle,|10\rangle\}\),
\[
X\otimes X + Y\otimes Y
\;\widehat{=}\;
2
\begin{bmatrix}
0 & 1 \\
1 & 0
\end{bmatrix}.
\]

Then
\begin{align*}
e^{-i\beta (X X + Y Y)}
&=
\cos(2\beta) I
- i \sin(2\beta)
\begin{bmatrix}
0 & 1 \\
1 & 0
\end{bmatrix}.
\end{align*}

So the statevector update
\[
\lvert 01\rangle \mapsto \cos(2\beta)\lvert 01\rangle - i\sin(2\beta)\lvert 10\rangle
\]
is consistent.

My wrong intermediate note was
\[
e^{-i\beta(X X + Y Y)} = \cos(\beta)I - i \sin(\beta) \sigma_x
\]
which misses the factor \(2\).

\entry{Wed 15 Oct, 11:40pm}{The edge-list bug that made the feasibility test fail}

The failing test was basically:

\begin{lstlisting}
every sampled bitstring should have exactly budget ones
\end{lstlisting}

and it was false. I was getting 3-ones strings from a budget-2 setup.

It was not the Hamiltonian sign. It was not the phase convention. It was not the Dicke state.
It was the edge pattern.

The actual XY edges in the code are alternating even edges, odd edges, and for even \(n\),
the wrap-around edge:

\begin{lstlisting}[language=Python]
even_edges = [(idx, idx + 1) for idx in range(0, n_assets - 1, 2)]
odd_edges = [(idx, idx + 1) for idx in range(1, n_assets - 1, 2)]
if n_assets % 2 == 0 and n_assets > 2:
    odd_edges.append((n_assets - 1, 0))
\end{lstlisting}

I had been reasoning as if the mixer acted on all pairs. Of course that leaks the intuition.
The whole constrained construction depends on the actual graph of pair exchanges.

\section*{Week 3 --- product-\(X\) mixer, QUBO-to-Ising mapping, and the Qiskit half-angle convention}

\entry{Mon 20 Oct, morning}{Product-\(X\) mixer is simpler, but it does \emph{not} preserve feasibility}

For a single qubit,
\[
e^{-i\beta X} = \cos(\beta)I - i\sin(\beta)X.
\]

So on the computational basis:
\begin{align*}
\lvert 0 \rangle &\mapsto \cos(\beta)\lvert 0 \rangle - i\sin(\beta)\lvert 1 \rangle, \\
\lvert 1 \rangle &\mapsto -i\sin(\beta)\lvert 0 \rangle + \cos(\beta)\lvert 1 \rangle.
\end{align*}

For \(n\) qubits the product mixer is
\[
U_B(\beta)
=
\prod_{j=1}^{n} e^{-i\beta X_j}.
\]

This is easy to implement and easy to reason about, but it does not preserve Hamming weight.
That means feasibility must be handled by the objective through the penalty, not by the mixer
geometry itself. This is why the repo supports both \repo{xy} and \repo{product\_x} as
\repo{mixer\_type}, and why the benchmark can compare them honestly instead of pretending one is
always better.

\entry{Tue 21 Oct, after two false starts}{Why the circuit uses \texorpdfstring{\repo{rx(2*beta)}}{rx(2 beta)}}

Qiskit's gate convention is
\[
R_x(\theta) = e^{-i \theta X / 2}.
\]

So if I want
\[
e^{-i \beta X}
\]
I need to set
\[
\theta = 2\beta.
\]

Same problem for cost terms:
\[
R_z(\theta) = e^{-i\theta Z / 2},
\qquad
R_{ZZ}(\theta) = e^{-i\theta Z\otimes Z / 2}.
\]

Therefore the code in \repo{build\_qaoa\_circuit(...)} is correct when it writes
\[
\repo{qc.rz(2.0 * gamma * coeff, qubit)}
\]
and
\[
\repo{qc.rzz(2.0 * gamma * coeff, q0, q1)}.
\]

My first attempt used \(\gamma\) and \(\beta\) directly and all the phases came out too small by
exactly a factor of \(2\). Annoying bug, but also a useful one because now I actually remember the
gate convention instead of memorising a line of code.

\entry{Thu 23 Oct, evening}{Deriving \repo{\_qubo\_to\_ising(...)}}

This was one of the more satisfying derivations because the code looked strange until I wrote it
out carefully.

For a symmetric QUBO matrix \(Q\),
\[
x^\top Q x

=
\sum_i Q_{ii} x_i
+
2\sum_{i<j} Q_{ij} x_i x_j.
\]

Use
\[
x_i = \frac{1 - z_i}{2},
\qquad z_i \in \{-1,+1\}.
\]

Then
\[
x_i x_j
=
\frac{1}{4}(1 - z_i - z_j + z_i z_j).
\]

So the diagonal term contributes
\[
Q_{ii}x_i = \frac{Q_{ii}}{2} - \frac{Q_{ii}}{2} z_i,
\]
meaning linear coefficient \(-Q_{ii}/2\).

For off-diagonal terms,
\begin{align*}
2Q_{ij}x_i x_j
&=
\frac{Q_{ij}}{2}(1 - z_i - z_j + z_i z_j).
\end{align*}

Therefore:
\begin{itemize}[leftmargin=1.4em]
\item quadratic Ising coupling is \(Q_{ij}/2\),
\item each involved qubit gets an additional linear shift \(-Q_{ij}/2\),
\item constants can be dropped.
\end{itemize}

That is exactly what the code does:
\begin{lstlisting}[language=Python]
linear[i] -= qubo[i, i] / 2.0
coupling = coeff / 2.0
linear[i] -= coeff / 2.0
linear[j] -= coeff / 2.0
\end{lstlisting}

I had an entire page of wrong algebra first because I forgot the symmetric QUBO contributes
\(2Q_{ij}x_ix_j\), not \(Q_{ij}x_ix_j\). Again: factor of \(2\), again: my own fault.

\section*{Week 4 --- objective function, periodic bounds, and finally understanding CVaR properly}

\entry{Mon 27 Oct, 8:30am}{Angle bounds and why periodic wrap matters}

The parameter vector ordering in the repo is
\[
[\gamma_0,\beta_0,\gamma_1,\beta_1,\dots]
\]
with bounds
\[
\gamma_\ell \in [0,2\pi),
\qquad
\beta_\ell \in [0,\pi).
\]

That becomes
\[
\text{\repo{make_bounds(cfg)}} \in \mathbb{R}^{2 \times 2p}.
\]

The projection step is
\[
\text{\repo{project\_params}}(x)
=
(x-\ell) \bmod (u-\ell) + \ell
\]
when periodic wrap is enabled.

This seemed almost too trivial to write down, until I realised BO was going to use Euclidean
distance on angle vectors. Without periodic wrap, \(0\) and \(2\pi - \varepsilon\) look far apart
even though physically they are almost the same point.

\entry{Tue 28 Oct, late afternoon}{CVaR by hand, because I kept saying the word without really owning it}

The repository uses a lower-tail CVaR objective for a minimisation problem:
sort energies from low to high, then average the best \(\alpha\)-fraction.

Tiny example:
\[
(-2.1,45),\quad (-1.8,120),\quad (-0.9,67),\quad (0.5,24)
\]
with total shots \(256\).

For \(\alpha=0.25\),
\[
\text{target weight} = 0.25 \times 256 = 64.
\]

Take all \(45\) shots from \(-2.1\), then \(19\) shots from \(-1.8\):
\[
\text{CVaR}
=
\frac{45(-2.1)+19(-1.8)}{64}
=
\frac{-94.5-34.2}{64}
=
-2.0109375.
\]

I kept saying ``upper tail'' in my notes for two days.
Wrong.
For minimisation this is the lower tail.

\entry{Wed 29 Oct, same page, more careful}{Tail variance and the BO noise proxy}

The repository does not just compute the lower-tail mean. It also computes
\[
\mathrm{Var}[X \mid \text{tail}] = \mathbb{E}[X^2 \mid \text{tail}] - \mathbb{E}[X \mid \text{tail}]^2
\]
and then uses
\[
\text{observation\_noise\_variance}
=
\frac{\mathrm{Var}[X \mid \text{tail}]}{\text{tail weight}}
\]
as a proxy for surrogate noise in BO.

Continuing the example above:
\[
\mathbb{E}[X^2 \mid \text{tail}]
=
\frac{45(2.1^2)+19(1.8^2)}{64}
=
\frac{45(4.41)+19(3.24)}{64}
=
\frac{198.45+61.56}{64}
=
4.06265625.
\]

So
\[
\mathrm{Var}[X \mid \text{tail}]
=
4.06265625 - (-2.0109375)^2
\approx 0.018786.
\]

And the noise proxy becomes
\[
\frac{0.018786}{64} \approx 2.94\times 10^{-4}.
\]

This was one of the moments where the code started to feel thought-through rather than improvised.
The variance is not pretending to be exact measurement uncertainty; it is acting as a useful
signal-quality proxy for the surrogate.

\section*{Week 5 --- random search, SPSA, and the first round of ``why is BO losing to random?''}

\entry{Mon 3 Nov, 7:15am somehow}{Random search is a real baseline, not an insult}

The random-search path is almost embarrassingly simple:

\begin{lstlisting}[language=Python]
for _ in range(cfg.evaluation_budget):
    params = rng.uniform(lower, upper)
    stats = evaluate_objective(...)
\end{lstlisting}

I almost apologised for it in my notes, which was silly. If random search is competitive,
that is a real result. It means the fancy outer loop may not justify its cost.

\entry{Tue 4 Nov, 10:40pm}{SPSA schedule written out numerically}

The code uses
\[
a_k = \frac{a}{(k+A)^{0.602}},
\qquad
c_k = \frac{c}{k^{0.101}}
\]
with \(a=0.2\), \(c=0.1\), and
\[
A = \max(5, \lfloor \text{evaluation\_budget}/5 \rfloor).
\]

For \(\text{evaluation\_budget}=40\),
\[
A = \max(5, 8) = 8.
\]

Then
\begin{align*}
k=1:
&&
a_1 &= \frac{0.2}{9^{0.602}} \approx 0.0534,
&
c_1 &= 0.1.
\\
k=5:
&&
a_5 &= \frac{0.2}{13^{0.602}} \approx 0.0420,
&
c_5 &= \frac{0.1}{5^{0.101}} \approx 0.0852.
\\
k=20:
&&
a_{20} &= \frac{0.2}{28^{0.602}} \approx 0.0270,
&
c_{20} &= \frac{0.1}{20^{0.101}} \approx 0.0739.
\end{align*}

That was enough to make the schedule feel real instead of ornamental.

\entry{Thu 6 Nov, frustrated enough to write this in capitals}{BO lost to random again}

One of the early small-budget runs gave me:
\[
\text{BO} = -0.2376,
\qquad
\text{random} = -0.2492.
\]

I wrote ``this cannot be right'' in my notebook, which was not fair.
It can absolutely be right. The first thing BO does is spend evaluations on the initial design.

If \(\text{\repo{evaluation\_budget}}=10\) and \(\text{\repo{n\_init\_points}}=8\), then the GP only gets to
be a GP for two actual adaptive steps. That is not a deep indictment of BO. It is just a budget
accounting fact.

\section*{Week 6 --- Bayesian optimisation details, periodic features, trust regions, and one overengineered but useful safeguard}

\entry{Mon 10 Nov, after staring at angle geometry for too long}{Periodic feature map}

The feature map in the BO code is
\[
\phi(x)
=
\left[
\cos\!\left(2\pi\frac{x-\ell}{u-\ell}\right),
\sin\!\left(2\pi\frac{x-\ell}{u-\ell}\right)
\right].
\]

This is the right thing to do for periodic angles.

For example, with \(\gamma \in [0,2\pi)\),
\[
\gamma_1 = 0,
\qquad
\gamma_2 = 2\pi - 10^{-6}
\]
are nearly identical physically, but in raw Euclidean coordinates they look far apart.
Under the feature map:
\[
\phi(\gamma_1) \approx (1,0),
\qquad
\phi(\gamma_2) \approx (1,-10^{-6}),
\]
which is exactly the behaviour I want.

\entry{Tue 11 Nov, night}{UCB vs EI for minimisation}

The code supports UCB and log-EI style behaviour.
For the sklearn fallback, the minimisation score is effectively
\[
\text{UCB score}(x) = \mu(x) - 2\sigma(x)
\]
and the improvement-style branch uses
\[
\text{EI}(x)
=
(f^\star - \mu(x))\Phi(z) + \sigma(x)\phi(z),
\qquad
z = \frac{f^\star - \mu(x)}{\sigma(x)}.
\]

Then the code negates the EI-style quantity because it still chooses by \(\arg\min\) over scores.

I had to write this down because I kept mixing up whether lower score meant more promising or less.

\entry{Wed 12 Nov, marginally smug}{Why random embedding is actually defensible}

This is the feature I thought might be overkill:
if the angle dimension exceeds \repo{bo\_max\_gp\_dim}, the code chooses a random active subspace.

Since dimension is \(2p\), if \(p=10\) then
\[
d = 2p = 20 > 16 = \text{\repo{bo\_max\_gp\_dim}}.
\]

Then the GP runs in \(16\) active dimensions rather than all \(20\).
I do not think this is the most important contribution in the repo, but it is at least a sober
acknowledgement that GPs do not become magically easy just because the variables happen to be QAOA angles.

\section*{Week 7 --- noise, fast simulator realism, measurement mitigation, and ZNE}

\entry{Mon 17 Nov, 9:30am}{Estimating gate load for the fast noise model}

The fast simulator does not pretend to be a full hardware noise model. It estimates an effective
bit-flip probability from the expected one- and two-qubit gate load.

Per layer:
\[
\text{cost two-qubit count} = \frac{n(n-1)}{2},
\qquad
\text{cost one-qubit count} = n.
\]

For the mixer:
\begin{itemize}[leftmargin=1.4em]
\item product-\(X\): one-qubit mixer count \(= n\),
\item XY: two-qubit mixer count \(= 2 \times |\text{XY edges}|\) because the circuit applies both
\repo{rxx} and \repo{ryy}.
\end{itemize}

For \(n=8\), \(p=2\), XY mixer:
\[
\frac{n(n-1)}{2} = \frac{8\cdot 7}{2} = 28
\]
cost couplers per layer.

For even \(n=8\), the XY edge list has \(8\) edges, so mixer two-qubit load per layer is
\[
2 \times 8 = 16.
\]

Therefore total two-qubit load estimate is
\[
2 \times (28 + 16) = 88.
\]

This is not hardware truth. It is an explicit approximation. But at least it is stated and
parameterised.

\entry{Tue 18 Nov, afternoon}{Fast depolarising proxy calculation}

The approximate fast noise probability in the code is
\[
p_{\text{flip}}
=
p_{\text{meas}}
+ (\text{oneq per qubit}) \cdot p_{1q}
+ (\text{twoq touches per qubit}) \cdot p_{2q}
\]
for the depolarising path.

If \(n=8\), \(p=2\), XY mixer, and I use the default
\[
p_{\text{meas}} = 0.01,\qquad p_{1q}=0.001,\qquad p_{2q}=0.01,
\]
then with the load estimate above:
\[
\text{oneq total} = 2 \times 8 = 16
\Rightarrow \text{oneq per qubit} = 2.
 \]

\[
\text{twoq touches per qubit}
=
\frac{2 \times 88}{8} = 22.
\]

So
\[
p_{\text{flip}} \approx 0.01 + 2(0.001) + 22(0.01) = 0.232.
\]

That is big. Maybe too big for some regimes. But it made me see why the code caps the final
probability and why mitigation can matter a lot in the noisy path.

\entry{Tue 18 Nov, later, after staring at \repo{simulator.py}}{Why the fast noise calibration is intentionally crude}

I added a note for myself here because I know a reviewer will ask where the constants came from.
The fast simulator uses
\[
p_{\text{flip}} \leftarrow 0.85 \times p_{\text{flip}}
\]
when twirling is enabled. This is not a first-principles theorem. It is a calibration discount:
the code is trying to reflect the empirical idea that randomized compiling / twirling tends to
``smear'' coherent over-rotation into something closer to stochastic noise, which is often easier
for mitigation to handle, but the repo does \emph{not} claim a device-faithful derivation.

The actual implementation logic in \repo{src/portfolio\_qaoa\_bench/simulator.py} is:
\begin{align*}
p_{\text{flip}}
&=
\min\Big(
0.45,\;
p_{\text{meas}}
+ g_{1q}\,p_{1q}
+ g_{2q}\,p_{2q}
\Big) \\
&\text{and then multiply by }0.85\text{ if twirling is on.}
\end{align*}

So there are really three modelling choices bundled together:
\begin{itemize}[leftmargin=1.4em]
\item gate-count compression into a single effective bit-flip rate,
\item hard clipping at \(0.45\) so the proxy does not become nonsense,
\item the \(0.85\) twirling discount as a pragmatic heuristic.
\end{itemize}

My own conclusion: acceptable for a benchmark proxy, absolutely not acceptable as a claim about
hardware physics. In the thesis I need to say this plainly. The value of the proxy is that it is
cheap, monotone in circuit/gate load, and exposes the optimisation loop to realistic degradation.
Its limitation is that it cannot distinguish coherent error, crosstalk, leakage, or correlated
readout effects. So if BO survives this model, that is evidence of robustness to \emph{a proxy},
not proof of hardware advantage.

\entry{Wed 19 Nov, finally a clean argument}{Measurement mitigation should not build the dense matrix if it can avoid it}

For one qubit, the confusion matrix is
\[
M =
\begin{bmatrix}
1-\varepsilon & \varepsilon \\
\varepsilon & 1-\varepsilon
\end{bmatrix}.
\]

For \(n\) qubits,
\[
p_{\text{obs}} = M^{\otimes n} p_{\text{true}}.
\]

My first instinct was to build \(M^{\otimes n}\) explicitly.
That is the sort of idea that seems fine until one does the arithmetic.

For \(n=12\):
\[
2^{12} = 4096,
\qquad
4096^2 = 16{,}777{,}216.
\]

If stored as float64:
\[
16{,}777{,}216 \times 8 \approx 134 \text{ MB}.
\]

For \(n=13\):
\[
8192^2 = 67{,}108{,}864,
\qquad
67{,}108{,}864 \times 8 \approx 537 \text{ MB}.
\]

So the better idea --- which is what the repo now does --- is:
\begin{enumerate}[leftmargin=1.4em]
\item invert the \(2\times 2\) single-qubit matrix once,
\item reshape the observed probabilities into an \(n\)-way tensor,
\item apply \(M^{-1}\) along each axis with tensor contraction,
\item renormalise and round back to shot counts.
\end{enumerate}

This is exactly the \repo{tensordot + moveaxis} loop in \repo{simulator.py}.

\entry{Thu 20 Nov, 11:50pm}{ZNE and the Richardson weights}

The ZNE path uses noise factors \((1,3,5)\) and solves for weights
\((w_1,w_3,w_5)\) such that constants are preserved and linear/quadratic noise terms cancel:
\begin{align*}
w_1 + w_3 + w_5 &= 1, \\
1\cdot w_1 + 3\cdot w_3 + 5\cdot w_5 &= 0, \\
1\cdot w_1 + 9\cdot w_3 + 25\cdot w_5 &= 0.
\end{align*}

Subtracting the second equation from the third:
\[
6w_3 + 20w_5 = 0
\qquad \Rightarrow \qquad
w_3 = -\frac{10}{3}w_5.
\]

Then from the second equation:
\[
w_1 + 3\left(-\frac{10}{3}w_5\right) + 5w_5 = 0
\Rightarrow
w_1 = 5w_5.
\]

Plug into the first equation:
\[
5w_5 - \frac{10}{3}w_5 + w_5 = 1
\Rightarrow
\frac{8}{3} w_5 = 1
\Rightarrow
w_5 = \frac{3}{8}.
\]

So
\[
w_5 = 0.375,\qquad
w_3 = -1.25,\qquad
w_1 = 1.875.
\]

Negative weight looked wrong the first time I saw it.
It is not wrong. It is the point of the extrapolation.

\section*{Week 8 --- pipeline, timing, the missing transpose bug, and why cost accounting became part of the science}

\entry{Mon 24 Nov, 11:47pm}{The missing transpose bug}

Three hours.
Actually three hours.

The bug was that the covariance matrix was not being symmetrised before QUBO construction.
I had effectively left \(\Sigma\) untouched when what I needed was
\[
\Sigma \leftarrow \frac{\Sigma + \Sigma^\top}{2}.
\]

That line now lives in \repo{QuboFactory.build(...)}:

\begin{lstlisting}[language=Python]
sigma = np.asarray(sigma, dtype=float)
sigma = (sigma + sigma.T) / 2.0
\end{lstlisting}

The symptom was not a dramatic crash. Much worse. On some seeds the matrix drifted just enough
to change small eigen-structure details and then penalty calibration started behaving strangely.
This is exactly the sort of bug that makes me distrust myself because the code \emph{almost} works.

\entry{Tue 25 Nov, calmer}{Shot-equivalent cost multiplier}

The repository now treats mitigation and robustness features as part of the effective shot cost.

In the code:
\[
m
=
1
\times
\begin{cases}
1.25 & \text{if measurement mitigation}\\
1 & \text{otherwise}
\end{cases}
\times
\begin{cases}
r & \text{if twirling with } r \text{ randomisations}\\
1 & \text{otherwise}
\end{cases}
\times
\begin{cases}
\sum_f \lambda_f & \text{if ZNE}\\
1 & \text{otherwise}
\end{cases}
\times
\begin{cases}
1+0.2L & \text{if resilience level } L > 0\\
1 & \text{otherwise.}
\end{cases}
\]

Example:
measurement mitigation on, twirling with \(r=4\), ZNE with \((1,3,5)\), resilience \(L=2\):
\[
m = 1.25 \times 4 \times 9 \times 1.4 = 63.
\]

That means \(256\) nominal shots become
\[
256 \times 63 = 16128
\]
effective shots in the accounting model.

This is one of the repo changes I am most convinced by. If BO wins only after burning through
enormous mitigation-adjusted cost, that should be visible.

\entry{Thu 27 Nov, afternoon}{Timing breakdown as an actual object, not a logging afterthought}

The repo now tracks
\[
T_{\text{total}}
=
T_{\text{classical}}
+ T_{\text{circuit}}
+ T_{\text{transpile}}
+ T_{\text{execution}}
+ T_{\text{mitigation}}
+ T_{\text{queue}}.
\]

This became the \repo{TimingBreakdown} dataclass in \repo{results.py}.

I originally thought this was just engineering neatness.
It is not.
The research question is explicitly about whether sophisticated outer-loop optimisation is worth
its real hybrid-loop cost. If I do not track the cost honestly, the question collapses.

\section*{Week 9 --- the supervisor conversation I deserved, the classical baseline, and the naming cleanup}

\entry{Mon 1 Dec, immediately after supervision}{The classical baseline comment landed because it was correct}

I said ``scope'' when asked why there was no classical Markowitz baseline.
That was a bad answer.

The repo was comparing BO vs SPSA vs random for tuning QAOA, but without a classical portfolio
baseline the question ``when is QAOA worth it?'' was underdefined.

I wrote this down in all caps in the notebook:
\[
\text{IF THE CLASSICAL SOLVE WINS CLEANLY, THAT IS STILL A RESULT.}
\]

\entry{Tue 2 Dec, evening}{Exact enumeration vs SLSQP relaxation}

The baseline now implemented in \repo{run\_classical\_markowitz(...)} has two branches.

For small enough systems:
\[
\min_{x \in \{0,1\}^n,\; \sum_i x_i = B} x^\top Q_{\text{base}} x
\]
by exact enumeration over supports.

For larger systems:
\begin{align*}
\min_{x \in [0,1]^n} \quad & x^\top Q_{\text{base}} x \\
\text{s.t.}\quad & \sum_i x_i = B,
\end{align*}
\]
solved by SLSQP, then rounded to the top-\(B\) entries.

It is not a perfect apples-to-apples baseline, but it is honest and interpretable.

\entry{Wed 3 Dec, late}{Readable names or no thesis}

I finally renamed the opaque internal classes because I could not keep pretending it was harmless.
The mapping that ended up in the repo is:
\begin{align*}
\repo{\_OrbitWeave\_41} &\to \repo{BayesianSurrogate},\\
\repo{\_LatchMurk\_47} &\to \repo{SearchTraceAccumulator},\\
\repo{\_ShardMint\_44} &\to \repo{PortfolioQUBO},\\
\repo{\_NixVault\_71} &\to \repo{RunConfig}.
\end{align*}

I left compatibility aliases in place because a hard break would have been unnecessary churn.
But the primary types now say what they mean.

\entry{Thu 4 Dec}{Removing fake-complexity modules}

I removed \repo{secrets.py} and \repo{secure\_buffer.py}.
There was no real secret-bearing data path in this project and keeping those files made the repo
look more elaborate than it was.

This was not a security decision.
It was an honesty decision.

\section*{Week 10 --- fixing the \repo{hard\_budget} regime and making instance difficulty measurable}

\entry{Mon 8 Dec, 2pm}{The regime name should correspond to something in the math}

The original \repo{hard\_budget} regime bothered me because the budget field was not actually doing
enough work inside the data generator.

The current logic is
\[
\text{budget fraction} = \frac{B}{n},
\qquad
\text{middle pressure} = 1 - 2\left|0.5 - \frac{B}{n}\right|.
\]

Then
\[
\text{edge} = 0.0025 - 0.0010 \cdot \text{middle pressure}
\]
for the return split, and covariance scale is also increased near the middle.

This made sense to me after I wrote down the story:
when the budget is near half the assets, there are many feasible states and many near-feasible
alternatives. The landscape should feel more knife-edge there than at extreme budget fractions.

\entry{Tue 9 Dec, finally a better x-axis}{Constraint hardness}

The most useful small function I added is
\[
H(n,B) = \frac{\binom{n}{B}}{2^n}.
\]

This says what fraction of the full Hilbert space is feasible under the cardinality constraint.

For the actual \repo{asset\_scale} study in the repository, \(B=3\) and
\(n \in \{6,8,10,12\}\), so:
\[
\begin{array}{c c c c}
n & \binom{n}{3} & 2^n & H(n,3) \\
\hline
6  & 20  & 64   & 0.3125 \\
8  & 56  & 256  & 0.21875 \\
10 & 120 & 1024 & 0.1171875 \\
12 & 220 & 4096 & 0.0537109
\end{array}
\]

That table helped more than I expected.
It gives an immediate hypothesis:
random should do well when \(H\) is large, everyone should suffer when \(H\) is tiny, and BO
should have the most room to help in the middle.

\entry{Wed 10 Dec, evening}{Structural diagnostics, not just winner tables}

I added \repo{profile\_instance(...)} in \repo{landscape.py} to compute:
\begin{align*}
\text{condition number} &= \frac{\lambda_{\max}}{\lambda_{\min}} \quad \text{(over nonzero magnitudes)},\\
\text{spectral gap} &= E^{\text{feasible}}_1 - E^{\text{feasible}}_0,\\
\text{energy gap to infeasible} &= \min E_{\text{infeasible}} - \min E_{\text{feasible}},\\
\text{participation ratio} &= \frac{1}{\text{ground-state degeneracy}},\\
\text{constraint hardness} &= \frac{\binom{n}{B}}{2^n}.
\end{align*}

This is not a proof that geometry predicts optimiser success.
It is a refusal to treat every QUBO instance as if it were structurally identical.

\section*{Week 11 --- results start to say something coherent}

\entry{Mon 15 Dec, 10:30pm}{The asset-scale study actually tells a story}

Using the committed results in \repo{results/multi\_regime\_suite/}, the aggregated mean best raw
objective for the asset-scale sweep is:
\[
\begin{array}{c c c c c}
n & H(n,3) & \text{Random} & \text{BO} & \text{SPSA} \\
\hline
6  & 0.3125    & -0.3880 & -0.3864 & -0.3678 \\
8  & 0.21875   & -0.3018 & -0.3146 & -0.3101 \\
10 & 0.1171875 & -0.2886 & -0.2969 & -0.2964 \\
12 & 0.0537109 & -0.2492 & -0.2376 & -0.2307
\end{array}
\]

This is exactly the shape I hoped would emerge:
BO is best in the middle, random is hard to beat when the feasible fraction is larger, and BO
falls back once the valid signal gets too sparse.

\entry{Tue 16 Dec, follow-up}{The \repo{hard\_budget} row is more subtle than I expected}

For \repo{regime = hard\_budget}, the aggregated committed results are approximately:
\[
\begin{array}{c c}
\text{method} & \text{mean best raw objective} \\
\hline
\text{classical} & -0.3140 \\
\text{random} & -0.3097 \\
\text{BO} & -0.3096 \\
\text{SPSA} & -0.3090
\end{array}
\]

The classical method wins, but the three QAOA-side methods are very close to one another.

That is interesting for a different reason than I expected.
It means the separation is not primarily in the final best feasible point. It is more about the
outer-loop cost, the raw CVaR-like objective, and how reliably the methods get there.

\entry{Wed 17 Dec, reading instead of running}{Barren plateau note}

I finally wrote the limitation down properly:
if gradient variance decays with \(n\), then even a well-engineered outer loop eventually hits a
signal-quality wall.

The repository already stores \repo{variance\_history}, so the limitation is not hypothetical.
It is at least measurable.

I am not claiming I proved a barren plateau.
I am saying that if I do not even mention the possibility for \(n=12\) and layered QAOA, someone
else will mention it for me, and they will be right to do so.

\section*{Week 12 --- statistics, regret, and finally being explicit when the exact reference is missing}

\entry{Mon 22 Dec, 8:15am}{Bootstrap confidence intervals}

The report aggregation computes bootstrap confidence intervals on means.
The logic is:
\begin{enumerate}[leftmargin=1.4em]
\item take the observed metric values across seeds,
\item resample with replacement,
\item compute the mean for each resample,
\item use the \(2.5\%\) and \(97.5\%\) quantiles as the interval.
\end{enumerate}

In symbols, if the resampled means are \(m^{(1)},\dots,m^{(R)}\), then
\[
\text{CI}_{95\%} = \left[q_{0.025}(m),\; q_{0.975}(m)\right].
\]

The repository also has a native acceleration path for this in \repo{native.py}, which I did not
expect to care about at first, but it made the report generation feel much less sluggish.

\entry{Tue 23 Dec, annoyed}{Why the Mann--Whitney \(p\)-values were all \(1.0\)}

I thought I had broken the stats code.
I had not.

The problem is that the current significance test is on \repo{best\_valid\_energy} across only
three seeds, and many runs land on exactly the same best feasible portfolio.

When the sample sets look like ties on top of ties, of course the test does not separate them.
That is not a bug.
That is a warning that the metric is not rich enough for the question I am asking.

\entry{Wed 24 Dec, this one mattered}{Approximation gap should not exist without a reference}

This became one of the repo changes I trust most:
\[
\text{approx gap}
=
\begin{cases}
\text{best valid} - \text{exact feasible reference}, & \text{if reference exists},\\
\text{NA}, & \text{otherwise.}
\end{cases}
\]

Not \(0\).
Not an empirical best-so-far.
Not a quiet fallback.
Just \(\text{NA}\).

This is now encoded via \repo{exact\_reference\_max\_assets} and
\repo{exact\_feasible\_energy = None}.

\entry{Fri 26 Dec, afternoon}{Regret plots are only useful if I remember what regret means}

The reporting code computes simple and cumulative regret:
\[
r_t = f_t - f^\star,
\qquad
R_t = \sum_{s=1}^{t} r_s.
\]

Here \(f_t\) is the best raw objective found up to evaluation \(t\), and \(f^\star\) is the
chosen reference level.

This is such a basic formula that I am annoyed I needed to write it down, but the plot is much
more interpretable once I do.

\section*{Week 13 --- writing the thesis around the code instead of pretending the code is the thesis}

\entry{Mon 5 Jan, abstract pain}{The framing kept trying to drift toward hype}

Bad version: ``This thesis demonstrates quantum advantage in portfolio optimisation.''

Less bad version:
\[
\text{This thesis presents a characterisation study of constrained QAOA portfolio optimisation}
\]
under shared hybrid-loop budgets with a classical baseline and structured regime sweeps.

That sentence is not beautiful, but it is at least attached to the actual repo.

\entry{Tue 6 Jan}{Related work finally went somewhere useful}

I added the related work file because the README was sounding strangely self-contained, which is
never a good sign in research writing.

The minimum set I needed was:
\begin{itemize}[leftmargin=1.4em]
\item Farhi et al.\ 2014 for QAOA,
\item Hadfield et al.\ for constrained mixers,
\item Barkoutsos et al.\ for CVaR-style variational objectives,
\item Egger et al.\ for portfolio-related QAOA / warm-start context,
\item Herman et al.\ for quantum finance survey context,
\item the BoTorch paper for the BO path I am actually using.
\end{itemize}

\entry{Thu 8 Jan, 10pm}{Architecture notes and the point of writing them down}

I finally wrote \repo{docs/ARCHITECTURE.md}.

The reason was not documentation virtue.
The reason was that once I had modular pieces
(\repo{data.py}, \repo{qubo.py}, \repo{simulator.py}, \repo{objective.py},
\repo{optimizers.py}, \repo{landscape.py}, \repo{reporting.py}, \repo{pipeline.py}),
I needed a short explanation of how the project actually fits together.

Otherwise I kept carrying the whole structure around in my head and dropping bits of it when I
tried to explain it quickly.

\section*{Week 14 --- CI, tests, monolith mirror, and the final repo becoming less embarrassing to show people}

\entry{Mon 12 Jan, 6:40pm}{Coverage and type-checking are not glamour work, but they changed how I wrote}

The CI now runs:

\begin{lstlisting}
- name: Coverage
  run: pytest --cov=portfolio_qaoa_bench --cov-report=term-missing -q

- name: Type check
  run: mypy src/portfolio_qaoa_bench --ignore-missing-imports
\end{lstlisting}

This made a real difference.
If the configuration objects, the payload schema, the reporting layer, and the pipeline all agree
under tests and mypy, then I am much less likely to write thesis prose based on a misremembered
implementation detail.

\entry{Tue 13 Jan, late}{Tests that actually correspond to research claims}

The repository test suite is not just smoke tests anymore. It now checks things like:
\begin{itemize}[leftmargin=1.4em]
\item exact reference cutoff behaviour,
\item regime-aware QUBO and data construction,
\item runtime accounting consistency,
\item optimizer fallback paths,
\item landscape and reporting logic,
\item monolith sync for the single-file mirror.
\end{itemize}

By the end of this pass the local suite was
\[
79 \text{ passed},\quad 1 \text{ skipped}.
\]

That is not a theorem, obviously.
But it is the first time the repo has felt stable enough that I trust the numbers coming out of it.

\entry{Wed 14 Jan}{The monolith request I did not expect}

I ended up adding \repo{monolith-full/portfolio\_qaoa\_bench\_monolith.py} plus a generator
script and sync check.

This is not a scientific contribution.
It is a repository usability contribution.

Still, it belongs in the final repo state because it answers a real maintenance problem:
if someone wants one file to inspect, there is now one generated source mirror, and CI checks that
it stays in sync with the modular source tree.

\section*{Week 15 --- reporting metrics, Overleaf packaging, and the things I should have measured earlier}

\entry{Mon 19 Jan, 9:05pm}{The cost-justification question needed cost-normalised metrics, obviously}

This should have been in the repo earlier.
The whole research framing is:
\[
\text{when does a more sophisticated optimiser justify its cost?}
\]
and somehow I had not made the cost-normalised scalars first-class outputs.

They are now explicit in \repo{SearchTrace.as\_summary()}:
\[
\text{best valid energy per shot}
=
\frac{E^{\text{best}}_{\text{valid}}}{S_{\text{effective}}},
\qquad
\text{best valid energy per USD}
=
\frac{E^{\text{best}}_{\text{valid}}}{C_{\text{QPU}}}.
\]

Here
\[
S_{\text{effective}}
=
\sum_{t=1}^{T} \text{effective shots at evaluation } t,
\qquad
C_{\text{QPU}}
=
\sum_{t=1}^{T} \text{estimated QPU cost at evaluation } t.
\]

I like these metrics because they force the discussion back to the actual hybrid loop.
If BO finds a slightly lower energy but only after much larger billed cost, I no longer need to
hand-wave that trade-off in the thesis prose.
It is just there in the exported summary.

\entry{Tue 20 Jan, after dinner}{Approximation ratio and regret AUC are better summary scalars than raw gap alone}

The approximation gap
\[
\Delta = E^{\text{best}}_{\text{valid}} - E^\star_{\text{exact}}
\]
is fine inside one regime, but not ideal across regimes whose energy scales are different.

So I added the ratio
\[
\rho
=
\frac{E^{\text{best}}_{\text{valid}}}{E^\star_{\text{exact}}}
\]
when the exact feasible reference exists.

I also finally promoted the trajectory summary I kept informally talking about:
the area under the regret curve.
With running best-valid history \(b_t\) and exact reference \(E^\star\),
\[
r_t = b_t - E^\star,
\]
and the code now computes
\[
\text{AUC}_{\text{regret}}
\approx
\operatorname{trapz}(r_1,\dots,r_T).
\]

That is what I wanted all along:
not only ``who won at the end?''
but also ``who spent the evaluation budget well over the whole run?''

\entry{Wed 21 Jan, slightly annoyed with myself}{Budget violation should be a histogram, not a yes/no flag}

The old summary told me feasibility hit rate and valid ratio.
Useful, but incomplete.

What I actually wanted to know was:
if a run is infeasible, how infeasible is it?

So each evaluation record now stores
\[
\text{violation counts} = \{0:n_0, 1:n_1, 2:n_2,\dots\},
\]
where violation level \(k\) means
\[
\left| \sum_i x_i - B \right| = k.
\]

This matters more than I expected for the \repo{hard\_budget} style cases.
Being off by one is a very different failure mode from spraying probability mass all over the
wrong Hamming sectors.

I also added the best-feasible Sharpe ratio:
\[
\text{Sharpe}(x)
=
\frac{\mu^\top x}{\sqrt{x^\top \Sigma x}}.
\]

That closes a loop I had left open for too long.
An energy improvement is nice, but if I cannot translate the selected portfolio back into a
recognisable finance metric then I am making the project harder to defend than it needs to be.

\entry{Thu 22 Jan, late, but this one felt worth doing}{Shot noise vs algorithmic variance}

The CVaR code already returns a tail variance term.
The new split I wrote down is:
\[
\sigma^2_{\text{shot}}
\approx
\frac{\sigma^2_{\text{tail}}}{N_{\text{shots}}},
\qquad
\sigma^2_{\text{alg}}
=
\max\!\left(\sigma^2_{\text{tail}} - \sigma^2_{\text{shot}},\, 0\right).
\]

It is not a perfect decomposition.
But it is better than mashing all uncertainty together and then pretending I know whether the fix
should be ``more shots'' or ``better optimisation''.

While I was in there I also started recording the ZNE pre-extrapolation objective and correction size:
\[
\Delta_{\text{ZNE}}
=
\left|E_{\text{mitigated}} - E_{\lambda=1}\right|.
\]

That one is important.
Without it, ZNE is just a checkbox.
With it, I can at least ask whether extrapolation materially changed the answer or mostly added
variance and cost.

\entry{Fri 23 Jan, cleaner statistics at last}{Pairwise tests should mean pairwise tests}

The earlier report only surfaced a \(p\)-value versus random search.
That was too narrow.

The aggregation now computes pairwise Mann--Whitney \(U\) tests across:
\[
\{\text{classical},\ \text{random},\ \text{SPSA},\ \text{BO}\}
\]
on best-valid energy distributions.

Then I apply the blunt but honest Bonferroni correction:
\[
p_{\text{adj}} = \min(1,\; m p),
\]
where \(m\) is the number of finite pairwise tests in that study slice.

Not elegant.
But respectable, and much harder to overstate.

\entry{Sat 24 Jan, admin work that was actually useful}{Making the journal uploadable to Overleaf}

I finally stopped treating the journal as a local oddity and packaged it like something I could
actually hand over.

The Overleaf-facing project now has a proper
\repo{main.tex}
entry point and sits in its own uploadable folder.
That sounds minor, but it changed how I wrote the final cleanup:
once I imagined another person compiling it in a browser, I became much less tolerant of half-finished
notes and untraceable references to ``the file somewhere in src''.

This is the same lesson as the rest of the repo, really.
If the artefact is meant to be read by someone else, structure is not cosmetic.
It is part of the research work.

\section*{Final pass --- what the final repo actually is}

\entry{Final week, written more slowly than the rest}{What changed from day 0 to the final repo}

At day 0 I had a vague mental model of QAOA, a portfolio objective I had not fully derived, and a
very dangerous willingness to speak confidently about things I had not yet checked.

By the final repo state I had:
\begin{itemize}[leftmargin=1.4em]
\item a QUBO construction I can derive by hand,
\item two mixer paths I can explain and contrast,
\item a statevector implementation, circuit construction path, and backend-aware execution path,
\item CVaR scoring with tail variance and a surrogate-noise proxy,
\item random, SPSA, BO, and classical Markowitz baselines under shared accounting,
\item explicit structure metrics (\(H(n,B)\), spectral gap, condition number, gap-to-infeasible),
\item precomputed suite results and report artifacts,
\item cleaned naming, honest limitations, stronger tests, and a generated monolith mirror.
\end{itemize}

The important thing is not that every idea worked.
Several did not.
The important thing is that the repo stopped hiding its own assumptions.

I know now why the XY factor of \(2\) is there.
I know why product-\(X\) needs penalty-based feasibility handling.
I know why the classical baseline belongs even when it makes the quantum story less flattering.
I know why \repo{exact\_reference\_max\_assets} and \repo{measurement\_mitigation\_max\_bits}
should be explicit config fields instead of implied limits.

Most of all: I know that the final repository is better because the mistakes forced the code to
become more explicit.

\vspace{1em}
\hrule
\vspace{0.8em}

\textbf{Build note.} This journal is now maintained as LaTeX source.
In this workspace there is no installed LaTeX engine (\repo{pdflatex}, \repo{xelatex}, and
\repo{lualatex} were unavailable when I checked), so I could not compile a PDF locally here.
The source is the canonical version.

\end{document}
